\documentclass[11pt]{article}
\usepackage[T1]{fontenc}
\usepackage{lmodern}
\usepackage[margin=1in]{geometry}
\usepackage{amsmath,amssymb,amsthm,mathtools}
\usepackage{hyperref}
\hypersetup{colorlinks=true,urlcolor=blue,linkcolor=blue}

\newtheorem{theorem}{Theorem}
\newtheorem{corollary}{Corollary}
\newtheorem{remark}{Remark}
\newcommand{\Om}{\Omega}
\newcommand{\E}{\mathcal E}
\newcommand{\F}{\mathcal F}
\newcommand{\dd}{\mathrm d}
\newcommand{\Real}{\operatorname{Re}}
\newcommand{\Lp}[1]{L^{#1}(\Om)}

\title{E13: The Band-Pass Amplitude Equation as a Negative Gradient Flow\\
\large A Complete Variational and Dissipation Proof}
\author{}
\date{October 26, 2026}

\begin{document}
\maketitle

\begin{abstract}
We prove that the complex band-pass amplitude equation
\(
\partial_t A=(r-a\Delta-b\Delta^2)A-\beta|A|^2A
\)
is the negative Wirtinger gradient flow of the associated Swift--Hohenberg
energy.  We compute the first variation, establish an exact energy--dissipation
identity, identify the normalization relative to the real $L^2$ gradient,
and recover the preferred Fourier shell.  The claim is made under explicit
periodic boundary and regularity assumptions; it does not imply global
well-posedness or convergence to a unique pattern.
\end{abstract}

\section{Setting and statement}
Let $\Om=\mathbb T^d$ be a flat periodic torus, with $1\leq d\leq3$,
and let $r\in\mathbb R$, $a,b>0$, and $\beta\in\mathbb R$.  Complex
$L^2$ is viewed as a real Hilbert space with inner product
\begin{equation}\label{eq:realinner}
 (u,v)_{\mathbb R}=\Real\int_\Om u\overline v\,\dd x.
\end{equation}
Define the real-valued energy for $A\in H^2(\Om;\mathbb C)$ by
\begin{equation}\label{eq:energy}
 \E[A]=\int_\Om\left(-r|A|^2-a|\nabla A|^2
             +b|\Delta A|^2+\frac\beta2|A|^4\right)\dd x.
\end{equation}
This energy is well defined for $d\leq3$ because $H^2(\Om)$ embeds
continuously into $L^\infty(\Om)$.

\begin{theorem}[E13 negative-gradient-flow identity]\label{thm:e13}
Suppose
\begin{equation}\label{eq:regularity}
 A\in C^1([0,T];H^2(\Om;\mathbb C))\cap
       C([0,T];H^4(\Om;\mathbb C))
\end{equation}
and $A$ solves, as an identity in $L^2(\Om)$,
\begin{equation}\label{eq:e13}
 \partial_t A=(r-a\Delta-b\Delta^2)A-\beta|A|^2A.
\end{equation}
Then the Wirtinger variational derivative (with the convention below) is
\begin{equation}\label{eq:wirtinger}
 \frac{\delta\E}{\delta\overline A}
   =-rA+a\Delta A+b\Delta^2 A+\beta|A|^2A,
\end{equation}
so that
\begin{equation}\label{eq:negative}
 \partial_t A=-\frac{\delta\E}{\delta\overline A}.
\end{equation}
Moreover, for every $t\in[0,T]$,
\begin{equation}\label{eq:dissip}
 \frac{\dd}{\dd t}\E[A(t)]
      =-2\|\partial_t A(t)\|_{L^2(\Om)}^2\leq0,
\end{equation}
and consequently
\begin{equation}\label{eq:balance}
 \E[A(T)]+2\int_0^T\|\partial_t A(t)\|_{L^2(\Om)}^2\,\dd t
       =\E[A(0)].
\end{equation}
Equality in \eqref{eq:dissip} at time $t$ holds if and only if
$\partial_t A(t)=0$ in $L^2(\Om)$, equivalently if $A(t)$ satisfies
\begin{equation}\label{eq:stationary}
 (r-a\Delta-b\Delta^2)A-\beta|A|^2A=0.
\end{equation}
\end{theorem}

\section{First variation: proof of the gradient formula}
\begin{proof}[Proof of Theorem~\ref{thm:e13}]
Fix $A\in H^4(\Om;\mathbb C)$ and a smooth periodic test function
$h\in C^\infty(\Om;\mathbb C)$. For real $\varepsilon$, put
$A_\varepsilon=A+\varepsilon h$. We use the convention
\begin{equation}\label{eq:variationconv}
 D\E[A](h)=2\Real\int_\Om
       \frac{\delta\E}{\delta\overline A}\,\overline h\,\dd x.
\end{equation}
This is the standard conjugate-variable/Wirtinger normalization for a
real-valued functional.

\paragraph{Quadratic amplitude term.}
Since
\[
 \left.\frac{\dd}{\dd\varepsilon}|A+\varepsilon h|^2
 \right|_{\varepsilon=0}=2\Real(A\overline h),
\]
the contribution of $-r\int|A|^2$ is
\begin{equation}\label{eq:var1}
 -2r\Real\int_\Om A\overline h\,\dd x.
\end{equation}

\paragraph{Gradient term.}
Differentiating $-a\int|\nabla A|^2$ gives
\[
 -2a\Real\int_\Om\nabla A\cdot\nabla\overline h\,\dd x.
\]
By periodic integration by parts,
$\int_\Om\nabla A\cdot\nabla\overline h
=-\int_\Om(\Delta A)\overline h$. Thus this contribution is
\begin{equation}\label{eq:var2}
 2a\Real\int_\Om(\Delta A)\overline h\,\dd x.
\end{equation}

\paragraph{Fourth-order term.}
Differentiating $b\int|\Delta A|^2$ yields
\[
 2b\Real\int_\Om(\Delta A)(\Delta\overline h)\,\dd x.
\]
Two periodic integrations by parts, or self-adjointness of $\Delta$,
give $\int(\Delta A)(\Delta\overline h)
=\int(\Delta^2 A)\overline h$. Therefore the contribution is
\begin{equation}\label{eq:var3}
 2b\Real\int_\Om(\Delta^2 A)\overline h\,\dd x.
\end{equation}

\paragraph{Quartic term.}
The pointwise identity
\[
 \left.\frac{\dd}{\dd\varepsilon}|A+\varepsilon h|^4
 \right|_{\varepsilon=0}
   =4|A|^2\Real(A\overline h)
\]
shows that the contribution of $(\beta/2)\int|A|^4$ is
\begin{equation}\label{eq:var4}
 2\beta\Real\int_\Om|A|^2A\overline h\,\dd x.
\end{equation}

\paragraph{Combination.}
Adding \eqref{eq:var1}--\eqref{eq:var4} gives
\begin{equation}\label{eq:firstvariation}
 D\E[A](h)=2\Real\int_\Om
 (-rA+a\Delta A+b\Delta^2A+\beta|A|^2A)\overline h\,\dd x.
\end{equation}
By density, \eqref{eq:firstvariation} holds for all $h\in H^2(\Om)$.
Comparison with \eqref{eq:variationconv} proves \eqref{eq:wirtinger}.
Taking its negative gives exactly the right-hand side of
\eqref{eq:e13}, establishing \eqref{eq:negative}.

\paragraph{Energy dissipation.}
The regularity \eqref{eq:regularity} justifies the $H^2$ chain rule:
\[
 \frac{\dd}{\dd t}\E[A(t)]=D\E[A(t)](\partial_t A(t)).
\]
Substituting $h=\partial_t A$ into \eqref{eq:firstvariation} and then
using \eqref{eq:negative},
\begin{align*}
 \frac{\dd}{\dd t}\E[A(t)]
 &=2\Real\int_\Om
     \frac{\delta\E}{\delta\overline A}
        \overline{\partial_t A}\,\dd x\\
 &=-2\Real\int_\Om
       \partial_t A\,\overline{\partial_t A}\,\dd x\\
 &=-2\|\partial_t A\|_{L^2(\Om)}^2.
\end{align*}
This proves \eqref{eq:dissip}. Integrating in time proves
\eqref{eq:balance}. Finally, the integrand is a squared norm, so it
vanishes exactly when $\partial_t A=0$, which is equivalent to
\eqref{eq:stationary} via \eqref{eq:e13}.
\end{proof}

\section{Normalization of the real $L^2$ gradient}
Because the real $L^2$ gradient is defined by
\[
 D\E[A](h)=(\nabla_{\mathbb R}\E[A],h)_{\mathbb R},
\]
comparison with \eqref{eq:firstvariation} implies
\begin{equation}\label{eq:gradientfactor}
 \nabla_{\mathbb R}\E[A]
   =2\frac{\delta\E}{\delta\overline A}.
\end{equation}
Consequently the equation can also be written
\[
 \partial_t A=-\tfrac12\nabla_{\mathbb R}\E[A].
\]
If $\F[A]=\E[A]/2$, then E13 is a literal negative real $L^2$
gradient flow:
\begin{equation}\label{eq:realflow}
 \partial_t A=-\nabla_{\mathbb R}\F[A],\qquad
 \frac{\dd}{\dd t}\F[A(t)]
     =-\|\partial_t A(t)\|_{L^2(\Om)}^2.
\end{equation}
The factor of two is a convention and does not change the dynamics.

\section{Centered Swift--Hohenberg representation}
Define
\begin{equation}\label{eq:kmu}
 k_\star^2=\frac{a}{2b},\qquad
 \mu=r+\frac{a^2}{4b}.
\end{equation}
Direct expansion shows
\[
 r-a\Delta-b\Delta^2=\mu-b(\Delta+k_\star^2)^2,
\]
so the same equation takes the centered form
\begin{equation}\label{eq:centeredflow}
 \partial_t A=\mu A-b(\Delta+k_\star^2)^2A-\beta|A|^2A.
\end{equation}
Consider
\begin{equation}\label{eq:centeredenergy}
 \E_{\mathrm{SH}}[A]
 =\int_\Om\left[-\mu|A|^2
      +b|(\Delta+k_\star^2)A|^2
      +\frac\beta2|A|^4\right]\dd x.
\end{equation}
Periodic integration by parts gives
\begin{align*}
 \|(\Delta+k_\star^2)A\|_2^2
 &=\|\Delta A\|_2^2
   +2k_\star^2\Real\int_\Om(\Delta A)\overline A\,\dd x
   +k_\star^4\|A\|_2^2\\
 &=\|\Delta A\|_2^2-2k_\star^2\|\nabla A\|_2^2
   +k_\star^4\|A\|_2^2.
\end{align*}
Since $2bk_\star^2=a$ and $bk_\star^4-\mu=-r$, it follows that
\begin{equation}\label{eq:equalenergy}
 \E_{\mathrm{SH}}[A]=\E[A]
 \qquad(A\in H^2(\Om;\mathbb C)).
\end{equation}
Thus the centered and uncentered energy descriptions are identical,
not merely equivalent up to an additive constant.

\section{Preferred Fourier shell}
Linearizing \eqref{eq:e13} at $A=0$ and substituting a Fourier
mode $e^{ik\cdot x}$ gives
\[
 \partial_t\widehat A_k=\lambda(k)\widehat A_k,
 \qquad
 \lambda(k)=r+a|k|^2-b|k|^4.
\]
Completion of the square using \eqref{eq:kmu} yields
\begin{equation}\label{eq:dispersion}
 \lambda(k)=\mu-b(|k|^2-k_\star^2)^2.
\end{equation}
In the continuum variable $q=|k|^2\geq0$, the function
$\lambda(q)=r+aq-bq^2$ has
\[
 \lambda'(q)=a-2bq,\qquad\lambda''(q)=-2b<0.
\]
Its unique maximizing radius is therefore $|k|=k_\star>0$.
When $\mu>0$, the linearly growing modes satisfy
\begin{equation}\label{eq:unstableband}
 \bigl||k|^2-k_\star^2\bigr|<\sqrt{\mu/b}.
\end{equation}
On the torus, only discrete lattice wavenumbers occur; the actual
fastest available modes maximize \eqref{eq:dispersion} over that lattice.
The nonlinear equation does not in general preserve a finite Fourier
shell exactly.

\section{Lower bound when $\beta>0$}
\begin{corollary}\label{cor:bound}
If $\beta>0$ and $\Om$ has finite volume, then
\begin{equation}\label{eq:lowerbound}
 \E[A]\geq-\frac{|\Om|}{2\beta}(\max\{\mu,0\})^2.
\end{equation}
If a solution satisfying the theorem's regularity assumptions exists
globally forward in time, then
\begin{equation}\label{eq:integrateddissipation}
 \int_0^\infty\|\partial_t A(t)\|_2^2\,\dd t<\infty.
\end{equation}
\end{corollary}
\begin{proof}
For $s\geq0$, completing the square (or minimizing the quadratic on
$[0,\infty)$) gives
\[
 \frac\beta2 s^2-\mu s
 \geq-\frac{(\max\{\mu,0\})^2}{2\beta}.
\]
Take $s=|A(x)|^2$ and integrate in \eqref{eq:centeredenergy}.
The remaining fourth-order square is nonnegative, proving
\eqref{eq:lowerbound}.  Integrate \eqref{eq:dissip} from $0$ to $T$,
apply the lower bound, and let $T\to\infty$ to obtain
\eqref{eq:integrateddissipation}.
\end{proof}

\begin{remark}[Scope]
The algebraic first-variation identity does not require $\beta>0$;
positivity is used only for the lower-bound corollary.  The proof
establishes monotone energy descent for regular solutions, but does
not by itself prove global well-posedness, convergence to a unique
steady state, nonlinear Fourier-shell invariance, or derivation of the
amplitude equation from a more fundamental physical model.  On other
domains the same computation holds when boundary conditions eliminate
all boundary terms and give the required self-adjoint operator.
\end{remark}

\section*{Conclusion}
Under the hypotheses of Theorem~\ref{thm:e13},
\begin{equation*}
 \boxed{\displaystyle
 \partial_t A=-\frac{\delta\E}{\delta\overline A},\qquad
 \frac{\dd}{\dd t}\E[A(t)]=-2\|\partial_t A(t)\|_2^2\leq0.}
\end{equation*}
This is the complete E13 negative-gradient-flow and Lyapunov-dissipation
identity for the specified periodic Swift--Hohenberg amplitude model.

\paragraph{Equation registry.}
\url{https://github.com/rickyjreyes/geometry_of_resonance/blob/main/WCT_FULL_EQUATION_LIST_CORRECTED.md}

\end{document}